\section{Certified programs and persistent tape equations}\label{sec:outer-programs}

We turn local code constraints into certificates for computations on the
tape. This section supplies local reversible filling
(Lemma~\ref{lem:outer-fill}) and a persistent tuple of tape equations
(Lemma~\ref{lem:outer-wiring}); Section~\ref{sec:outer-representation}
will assemble them into the evolving representation.

We retain the outer notation of Sections~\ref{sec:macrograph} and
\ref{sec:outer-codes}: $h=2^\ell$, $M_o=2^r$, the odd-characteristic
field $J$, the group $\Gamma=\mathbb F_2^r\rtimes H$, and coordinates
$k=(g,u)=(ux,u)$.  In particular $r=O(\log N)$,
$|H|\le r^9$, and a field element has $O(\log N)$ bits.
All constants and polynomial exponents in this section are fixed
independently of $N$.  The local star degree $d$ is a constant.
We pack the four parts $(s,t)\in\{0,1\}^2$ at one $k$ into one symbol;
``part distance'' instead normalizes over the four parts separately.
These distances differ by at most a factor of four.

There are two kinds of storage.  A \emph{full slot} has both axis
constraints and duplicated systematic payload, and thus equals
$E_{\alpha,\beta}v$ for a unique payload $v:\Gamma\to U$.
An \emph{A-only scratch slot} has the same local A constraints and
A-edge agreements, but its other labels and entry coordinates are
merely vector coordinates: no B condition is imposed.  Keeping all
intermediate certificates, rather than just the values being computed,
will let us certify a computation at every point of a reversible schedule.

\subsection{Full operations and a fixed query list}

\begin{definition}[Certified full operation]\label{op:full-operation}
A full operation specifies parent full slots, a target encoder
$E_{\mathrm{out}}$, and a local expression $X^{st}_k(i,j)$ in the parents.
For arbitrary exact duplicated full parents, $X$ must have duplicated
systematic payload $v$ and satisfy an allowed A-axis constraint system.
In addition, either its A system is that of the target, or $X$ satisfies
the target's B system.  The required target is $E_{\mathrm{out}}v$.
Scalar coefficients in linear combinations are identical at every
index.  They may depend on an identical clock and on previously
certified, completed public address data.  The specified covariant
primitive $W(k)\mapsto W(k)\rho_e(k)$ below may additionally use its
explicit index-dependent public matrix; it must satisfy the same
axis-constraint requirements on $X$.
\end{definition}

In the first case compare $X$ and the target on all systematic B columns
at $t=0$, retaining every entry in each such column at both A parts.
In the second case compare all systematic A rows at $s=0$, retaining
their entire rows at both B parts.  The one-axis uniqueness and distance
in Theorem~\ref{thm:outer-codes} give a constant $\gamma_{\mathrm{op}}>0$
with the following alternative: the required payload identity is true
and all comparisons hold, or these comparisons fail on at least
$\gamma_{\mathrm{op}}$ of the centers.  Indeed, equality of payloads
fixes each compared column or row by the shared one-axis extension.
If a payload differs, at least one compared column or row is a nonzero
codeword difference for that axis, so its vertex support has constant
density.  The number of labels and parts is fixed.  A comparison of
isolated systematic cells would not have this implication.

We use the following full operations, splitting sums into bounded-arity
stages whenever necessary.
\begin{enumerate}[label=(\roman*)]
\item Scalar linear combinations and right pulls
$R_bF(k)=F(kb)$ share the scalar A system.  Elementary pulls are the
identity, the three first-bit translations $(e_a,1)$, and one-step
Singer-cycle rotations of each block and their inverses.
\item A matrix payload can be reencoded while changing only one axis
type, using the other as the shared axis.  Constant linear maps and
contractions of fully untwisted entry alphabets also give full operations.
\item To multiply a payload by a one-bit Walsh character, we must also
certify the resulting full encoding. An index-dependent scalar
multiplier need not preserve the scalar edge equations. The matrix
representation supplies a shared-axis route through the twisted types.
Starting from an untwisted scalar slot $W$, form
\[
             X^{st}_k(i,j)=W^{st}_k(i,j)\rho_e(k),
\]
where $e$ is a distinguished first bit of a block and $\rho_e$ is the
representation of Section~\ref{sec:outer-codes}.  Multiplicativity gives
\[
 X^{1t}_{a_i k}(i,j)=\rho_e(a_i)X^{0t}_k(i,j),\qquad
 X^{s1}_{k b_j}(i,j)=X^{s0}_k(i,j)\rho_e(b_j).
\]
Scalar local checks continue to hold, since the same entry matrix
multiplies every local scalar.  Reencode first with twisted A and
untwisted B, sharing A, and then with both axes untwisted, sharing B.
The column of $\rho_e(g,u)$ indexed by $1_H$ has just one nonzero entry,
in row $u$, equal to
$(-1)^{(ue)\cdot g}=(-1)^{e\cdot x}$.
Summing this entry-alphabet column therefore produces the full scalar
encoding of $(-1)^{e\cdot x}v_W(k)$.  The contraction is on entry
coordinates, not on outer star labels.
\end{enumerate}
Every intermediate matrix slot in this construction has a specified
homogeneous linear payload recipe.  An arbitrary multiplier depending
on one bit, with values $a_0,a_1\in J$, is
$\frac12(a_0+a_1)I+\frac12(a_0-a_1)M_e$, where
$M_ef(x,u)=(-1)^{e\cdot x}f(x,u)$ is the corresponding modulation.
The divisions by two are valid because $J$ has odd characteristic.

Right rotations move these operations to any bit in a block.  For a
block rotation $v$ the pull convention gives
\[
 R_{(0,v)}R_{(e,1)}R_{(0,v^{-1})}=R_{(ve,1)}.
\]
Conjugating the modulation in the same way gives the modulation of bit
$ve$.  A power of a rotation is implemented as a sequence of single
rotations, returning to the original layout after the operation.
Thus an arbitrary bit operation costs $O(r)$ elementary full stages,
without introducing new query maps.

Fix once and for all a permutation list $\Pi$ containing identity,
the required forward and inverse A and B neighbor maps, left pulls by
a fixed generating list of $H$, the elementary right pulls above,
and right pulls by a fixed generating list of $H$ and its inverses.
Include any bounded compositions used in a local expression.
Its cardinality $q=|\Pi|$ is independent of $N$.
All constraints of all slots are checked together at a center using
the symbols at $\Pi(k)$.  A growing number of slots consequently
increases local computation and symbol size, but not the number of
queried positions.  In particular a single false exact full operation
retains the gap $\gamma_{\mathrm{op}}$, rather than having its gap
divided by the program length.

\subsection{Scalar transforms with certified stages}

\begin{lemma}[Scalar program library]\label{op:scalar-programs}
On scalar tables indexed by $x\in\mathbb F_2^r$, the full operations
above give homogeneous linear programs of $\poly(r)$ stages for:
bit-conjunction and interval masks; controlled bit flips and bit
permutations; integer cyclic shifts; discrete Fourier transforms and
their inverses of lengths $M_o$ and $h$; the Walsh transform; and cyclic
coefficient multiplication by any systematic Lagrange polynomial
$L_j$, $0\le j<h$.  The last operation is uniform also when $j$ is
supplied by already identical address bits.  All intermediate stages
are full slots with fixed encoder types and query maps in $\Pi$.
\end{lemma}

\begin{proof}
A bit-conjunction mask is a product of one-bit masks, applied in
sequence.  For the integer interval $[0,a)$ with $0\le a\le M_o$,
represent $a$ by $r+1$ public or already identical control bits.
Use the zero operator for $a=0$ and the identity for $a=M_o$.
For $0<a<M_o$, use the disjoint prefix rectangles obtained by
specifying the most significant differing bit: at a position where
$a$ has bit one, match its more significant bits and require the
current bit to be zero.  There are at most $r$ rectangles.
More explicitly, writing $a_r$ for the endpoint bit, one fixed
padded program applies the mask
\[
 a_r+(1-a_r)\sum_{i=0}^{r-1}a_i\,
        1[x_i=0]\prod_{j=i+1}^{r-1}1[x_j=a_j].
\]
All endpoint decisions and rectangle coefficients are identical at
every index; the program keeps all candidate stages and uses scalar
coefficient selection, so its topology and polynomial size do not
depend on $a$.  Differences of two such masks give arbitrary intervals.

Let $m$ be the indicator of a conjunction not involving bit $i$.
The pullback of the permutation that flips bit $i$ precisely when
$m=1$ is
\begin{equation}\label{op:controlled-flip}
                    f-mf+R_{(e_i,1)}(mf).
\end{equation}
To verify this, if $m(x)=0$ then $m(x+e_i)=0$, and the result is $f(x)$;
if $m(x)=1$ then $m(x+e_i)=1$, and the result is $f(x+e_i)$.
The independence of $m$ from the flipped bit is essential here.
Three controlled XOR operations exchange two index bits, so they
implement bit reversal and any prescribed bit permutation.

For an integer subtraction of $2^j$ modulo $2^r$, flip bit $i>j$
when the original bits $j,\ldots,i-1$ are all zero, and flip bit $j$
unconditionally.  Process the higher bits in descending order;
then none of the bits used by a condition has yet changed.
The addition rule uses ones instead of zeros.  Pullbacks reverse the
chronological composition of index permutations: thus the table
stages apply these borrow gates in ascending bit order, the reverse
of their descending index-permutation order.  This implements the
desired cyclic table shift.  A conjunction uses $O(r)$ bit
operations, so this procedure is polynomial in $r$.

For the Fourier transform we realize the radix-two Cooley--Tukey
decomposition~\cite[pp.~299--300]{CooleyTukey1965} by the certified full
operations above. Let $D=2^v\in\{h,M_o\}$ and let
$\xi=\zeta^{M_o/D}$ have order $D$.
The transform acts on the low $v$ integer bits, independently within
each block of fixed higher bits, and is
\begin{equation}\label{op:dft}
       \mathcal F_D(f)(y)=\sum_{z=0}^{D-1}\xi^{yz}f(z).
\end{equation}
First reverse these $v$ bits.  At level $j$, pair positions separated
by XOR of bit $j-1$ in each aligned window of length $2^j$.
Writing $a<2^{j-1}$ for the offset and $U,V$ for the lower and upper
entries, replace the pair by
\begin{equation}\label{op:butterfly}
          U+\xi^{aD/2^j}V,\qquad U-\xi^{aD/2^j}V.
\end{equation}
The factor is a product of one-bit factors, since if
$a=\sum_{i<j-1}2^ia_i$, then
$\xi^{aD/2^j}=\prod_{i<j-1}(\xi^{2^iD/2^j})^{a_i}$.
Half-window masks, the bit pull, these factors, and addition implement
the displayed rule.  The identity
\[
 \sum_{z=0}^{D-1}\xi^{yz}f(z)
 =\sum_{z=0}^{D/2-1}(\xi^2)^{yz}f(2z)
  +\xi^y\sum_{z=0}^{D/2-1}(\xi^2)^{yz}f(2z+1)
\]
proves the algorithm by induction, with bit reversal placing the
recursive even and odd subproblems in the required order.  More
explicitly, if $\operatorname{rev}_a$ reverses $a$ bits, after level
$j$ the entry at $b2^j+y$, $0\le y<2^j$, is
\[
 \sum_{z=0}^{2^j-1}(\xi^{D/2^j})^{yz}
 f\bigl(\operatorname{rev}_{v-j}(b)+2^{v-j}z\bigr).
\]
The level-$j$ butterfly combines the even and odd values of $z$ in
the level-$(j-1)$ formula, proving this invariant from the bit-reversed
input at level zero; at level $v$ it is \eqref{op:dft}.
Using $\xi^{-1}$ and multiplying by $D^{-1}$ gives the inverse:
$\sum_{y=0}^{D-1}\xi^{y(z-z')}=D\,1[z=z']$.
Without initial bit reversal, apply the sum/difference butterflies
directly to the input table, one coordinate bit at a time, and use
no twiddle factors.  This gives
\begin{equation}\label{op:walsh}
                   \mathcal W(f)(x)=\sum_y(-1)^{x\cdot y}f(y).
\end{equation}
They can be applied on all $r$ bits or just a specified block.
The Walsh inverse is $M_o^{-1}\mathcal W$ when all bits are used.

Finally set $\omega=\zeta^{M_o/h}$ and
\begin{equation}\label{op:lagrange}
 L_j(Z)=h^{-1}\sum_{a=0}^{h-1}\omega^{-ja}Z^a.
\end{equation}
The root sum just used gives $L_j(\omega^i)=1[i=j]$.
For a length-$M_o$ coefficient table $F$, cyclic multiplication by
$L_j$ has coefficient at $w$ equal to
\begin{equation}\label{op:lagrange-convolution}
 h^{-1}\omega^{-jw}\sum_{a=0}^{h-1}F_{w-a}\omega^{j(w-a)},
\end{equation}
with subscripts modulo $M_o$.  The powers are well defined under
cyclic indexing because $h\mid M_o$.  Each of the two diagonal
factors splits over the bits of $w$ and has coefficients computable
from $j$.  If $T_bf(w)=f(w-b\bmod M_o)$, the middle window sum is
\[
                    \prod_{i=0}^{\ell-1}(I+T_{2^i}).
\]
Every integer in $[0,h)$ has a unique binary expansion, proving this
identity.  The previously constructed shifts implement it.
When $\deg(F L_j)<M_o$, cyclic and ordinary multiplication coincide;
we impose this degree bound at every ordinary-product use.

For clarity, crude stage bounds suffice: a bit operation uses $O(r)$
elementary stages, a conjunction $O(r^2)$, a carry shift $O(r^3)$,
and the window product $O(r^4)$.  The Fourier, Walsh, and range-mask
constructions also fit a fixed polynomial bound.  Coefficient
calculation, layout descriptions and every elementary local operation
use polynomially many bit operations in $r$.  Padding unused steps
with zero or identity operations makes all lengths independent of
the particular control or address values.
\end{proof}

\subsection{Triangular A-only certificates}

An A-only program starts from A-coded public constants or a local
expression $X$ from Definition~\ref{op:full-operation}.  It forms
identical-coefficient linear combinations and right pulls of earlier
slots having the same A twist.  Right pulls commute with left
translations and constant entry twists, and therefore preserve A
constraints.  A fixed generating list of $H$ suffices to compute all
desired right-$H$ pulls through stored paths in its Cayley graph.
Each path has length at most $|H|-1$ after removing repetitions.
Arbitrary bit pulls use the rotation sequences already described.
The number of stored stages is polynomial in $r$, while all recurrence
queries still belong to $\Pi$.

At a fixed clock value, the initialized slots form triangular prefixes
and satisfy their recurrences; every other scratch slot is zero.
Filling one next stage evaluates its recurrence locally.  Erasing in
reverse order replaces that stage by zero while its parents are still
available.  Each step changes one designated slot or fixed family of
slots, with the same definition of every retained slot.

\begin{lemma}[Simultaneous scratch restoration]\label{op:scratch-restoration}
Suppose all scratch recurrences at a legal repeated clock are
triangular and their references, on arbitrary exact duplicated full
parents, satisfy their assigned A constraints.  Public address-focus
recurrences precede every recurrence using the resulting coefficients.
There exist constants $C_{\mathrm{sc}}$ and $\theta_{\mathrm{sc}}>0$,
independent of the number of stages, slots and entry coordinates,
with the following property.  If the packed parsing, clock, full-code,
duplication, scratch local-code, scratch-edge, recurrence and inactive-zero
tests reject on a fraction $\theta<\theta_{\mathrm{sc}}$ of centers,
then after one preprocessing of the parsing, clock and full slots,
all scratch can be replaced simultaneously by its exact triangular
reference at total symbol cost at most $C_{\mathrm{sc}}\theta$.
The references are taken on the corrected full parents; those parents
need not yet satisfy their advertised payload-operation identities.
\end{lemma}

\begin{proof}
Theorem~\ref{thm:outer-codes} first repairs parsing, repetition and
full codes at $O(\theta)$ symbol cost.  A repeated illegal clock is
excluded by its range test.  The full-code duplication gap then
excludes any failure of duplicated payload, for small enough $\theta$.
Replace all locally A-invalid scratch parts by zero.  This is one
further packed $O(\theta)$ change.  Call the resulting word $W$, and
\emph{keep this word fixed throughout the following induction}.

If preprocessing changed at most $C_{\mathrm{pre}}\theta$ packed
sites, any test using $\Pi$ can change its outcome only on their
$q$ inverse images.  After harmless fixed conversions of edge and
part normalizations there are constants $C_z,C_r$ such that
\begin{equation}\label{op:fixed-word-invariant}
 z_A(W)\le C_z\theta,\qquad
 \Pr[\text{some scratch recurrence or inactive-zero test fails on }W]
       \le C_r\theta.
\end{equation}
One can use constant multiples of $1+qC_{\mathrm{pre}}$ for both
constants.  These bounds concern the \emph{entire} packed scratch
tuple, so they contain no slot-count factor.

Order all initialized families topologically, putting public focus
first, and include the inactive-zero families with their zero
references.  Let $V_j$ be the exact reference of the first $j$ families
on the corrected full parents and legal clock, and let $p_j$ be the
joint part distance from the corresponding prefix of $W$ to $V_j$.
The empty prefix has $p_0=0$.  The full parents may have incorrect
advertised identities, but are A-coded; every allowed initial $X$
is A-coded on such parents.  Right pulls and identical coefficients
preserve this property.  Public focus begins with hardwired correct
constant encodings and the identical clock.  Its completed reference
supplies globally identical coefficients to subsequent recurrences.
Consequently every $V_j$ is exactly A-coded.

An error in the new family occurs only at a recurrence-failure center
or at a center reading an erroneous earlier part.  A joint part-error
set of density $p_j$ occupies at most $4p_j$ of the packed sites.
Taking a union over the $q$ maps, and retaining the old prefix, gives
\begin{equation}\label{op:prefix-provisional}
                 p_{j+1}\le C_r\theta+(1+4q)p_j.
\end{equation}
This remains true when a recurrence reads all earlier slots or types
at one of those sites: the error event is already joint.
Put $K=2C_z/\delta_1$ and choose $\theta_{\mathrm{sc}}$ so small that
\begin{equation}\label{op:scratch-smallness}
          (C_r+(1+4q)K)\theta_{\mathrm{sc}}<\delta_1/16.
\end{equation}
If $p_j\le K\theta$, the provisional estimate is in the small-error
range of the A near-code inequality of Theorem~\ref{thm:outer-codes}.
Apply that inequality to the \emph{same fixed word} $W$, restricted
to the enlarged tuple, and to $V_{j+1}$:
\[
 \delta_1p_{j+1}
 \le z_A(W|_{j+1})+4p_{j+1}^2+\lambda p_{j+1}.
\]
The conflict count of a prefix is at most that of the whole tuple.
Since $\lambda\le\delta_1/8$ and $p_{j+1}<\delta_1/16$, this implies
\[
                 p_{j+1}\le 2z_A(W)/\delta_1\le K\theta.
\]
Induction proves this bound for the complete scratch tuple.
It resets an \emph{estimate}, without making successive word repairs.
Only now replace the whole scratch tuple by its reference, at symbol
cost at most $4K\theta$.  The common scalar local checks act on joint
vectors, and the finitely many twist types act by a block-diagonal
invertible map; the near-code estimate is therefore valid in exactly
this joint metric.  Together with preprocessing this proves the lemma.
\end{proof}

\subsection{Mode-serial filling of a full slot}

The target must remain a full codeword while it is being filled, even
though its final value can depend on the parents at every index. We use
a sequence of frequency cutoffs that commute with the encoder. Each
increment is supported on one $H$-orbit, so its encoded contribution
can be computed from a polylogarithmic list of coefficients gathered
locally. A completed certificate for the old cutoff is retained until
the next cutoff and its increment have both been prepared.

Retaining the completed target while reversing its auxiliary computation
follows the reversible-cleanup principle of Bennett~\cite[pp.~525--526]{Bennett1973}.
The two-buffer schedule, local transitions, and certification guarantees
needed here are established in the following proof.

\begin{lemma}[Local reversible filling]\label{lem:outer-fill}
For every full operation in Definition~\ref{op:full-operation}, there
is a canonical reversible schedule from zero target and zero temporary
scratch to its required full target and zero temporary scratch.
Parents and all coefficient data remain fixed throughout.  The schedule
has $O(M_o\poly(r))$ steps, a $\poly(r)$-bit local symbol and local
algorithm, and an $O(\log M_o+\log\poly(r))$-bit microclock.
Its checks and both local transition directions use only $\Pi$.
At each clock location the target is a specified full encoding and
all scratch consists of specified initialized prefixes.  After
simultaneous scratch restoration, failure of the required target
identity has a constant center gap.  Each reference scratch slot
satisfies its A constraints, even when its full parents have not yet
been shown to satisfy their advertised payload identities.
\end{lemma}

\begin{proof}
Write the operation's local expression as $X$ and its duplicated
systematic payload as $v$.  For each label, part and entry coordinate
expand in the $x$ variable:
\[
 X(x,u)=\sum_{\nu\in\mathbb F_2^r}\widehat X_\nu(u)(-1)^{x\cdot\nu},
 \qquad
 \widehat X_\nu(u)=M_o^{-1}\sum_xX(x,u)(-1)^{x\cdot\nu}.
\]
For $0\le\tau\le M_o$ define a Fourier multiplier $P_\tau$ by
\begin{equation}\label{op:cutoff}
 w_\tau(\nu)=|H|^{-1}\sum_{a\in H}
                    1[\underline{a^{-1}\nu}<\tau].
\end{equation}
Thus $P_0=0$ and $P_{M_o}=I$.  The inverses of $|H|$ and $M_o$
exist in $J$.  The multipliers are constant on frequency orbits of
$H$; they are not asserted to be projections.

Here is the required commutation in detail.  A left pull by $(w,a)$
sends $(x,u)$ to $(x+u^{-1}a^{-1}w,au)$, so it preserves each
frequency, changing its coefficient by a known sign and an $H$-index
permutation.  A right pull sends $(x,u)$ to
$(a^{-1}(x+w),ua)$, carrying a frequency $\nu$ to $a\nu$, with
sign $(-1)^{w\cdot a\nu}$.  Hence $P_\tau$ commutes with both pulls.
The constant entry twists act on a different coordinate and also
commute.  Every parity operator and every systematic injection thus
intertwines $P_\tau$ on its domain and codomain.  If an invertible
parity operator $T$ satisfies $TP_\tau=P_\tau T$, then its inverse
does too, by multiplication by $T^{-1}$ on both sides.  Applying the
two one-axis extensions gives
\begin{equation}\label{op:cutoff-extension}
              P_\tau E_{\mathrm{out}}v
                      =E_{\mathrm{out}}(P_\tau v).
\end{equation}
No commutation with a preliminary index-dependent operation such as
$W(k)\mapsto W(k)\rho_e(k)$ is needed: $P_\tau$ is applied to its
already formed expression $X$.

The mode-bit selector
\begin{equation}\label{op:mode-selector}
                   \frac{I+(-1)^bR_{(e_i,1)}}2
\end{equation}
retains precisely those frequencies with $\nu_i=b$.
For each $a\in H$, decompose
$\underline{a^{-1}\nu}<\tau$ into at most $r$ disjoint prefix
rectangles and multiply their selectors.  Sum over rectangles and
average over $a$.  This gives an A-only program computing $P_\tau X$
with $\poly(r)$ stages.  Include all candidate rectangles and use
coefficient bits from the microclock, so topology and length are fixed.
The case $\tau=M_o$ uses the identity with the same padding.
All selectors use the elementary pull list via stored bit alignments.

The partial target at cutoff $\tau$ is
$Y_\tau=E_{\mathrm{out}}(P_\tau v)$.
Compare it with the completed cutoff program using the shared-axis
test of Definition~\ref{op:full-operation}.  In the shared-B case,
intermediate scratch is only A-coded.  The test's justification is
made \emph{after} Lemma~\ref{op:scratch-restoration}: the completed
reference is then exactly $P_\tau X$, which satisfies the target's
B constraints by the commutation above.  One may now use B distance
on its whole systematic A rows.  This argument neither assumes nor
tests B validity of unfinished scratch.  The shared-A case uses the
analogous columns.  Thus equality of the required payload suffices,
and its failure has the same constant gap at every cutoff.

It remains to compute a cutoff increment locally.  Let $a_\tau$ be
the bit vector for the integer $\tau<M_o$, and set $O=Ha_\tau$.
The increment multiplier is zero outside $O$ and, for $\nu\in O$, is
\begin{equation}\label{op:orbit-weight}
 (w_{\tau+1}-w_\tau)(\nu)
 =\frac{\#\{a\in H:aa_\tau=\nu\}}{|H|}
 =\frac1{|O|}.
\end{equation}
For each $\nu$ in the fixed-length list $(aa_\tau)_{a\in H}$,
apply all $r$ mode selectors to $X$.  Repetitions are harmless.
The resulting pure-mode slot is
$X_\nu(x,u)=(-1)^{x\cdot\nu}\widehat X_\nu(u)$.
Also store each right pull $R_{(0,b)}X_\nu$, $b\in H$, through
generator paths and their intermediate stages.  At the current
$k=(x,u)$ its systematic payload is
\begin{equation}\label{op:orbit-focus}
             (-1)^{(b^{-1}x)\cdot\nu}\widehat v_\nu(ub).
\end{equation}
For any desired $u'\in H$, take $b=u^{-1}u'$ and remove the displayed
known sign.  One symbol now contains every coefficient
$\widehat v_\nu(u')$ for $\nu\in O$ and every $u'\in H$.
These are read from locally stored pull slots, not from remote global
payload positions.

We can apply $E_{\mathrm{out}}$ to this coefficient list locally.
Remark~\ref{oc:orbit-invariance} supplies the inverse of each one-axis
parity system on the $O$-frequency subspace: both this subspace and its
complement are invariant, and the restricted domain and codomain have
equal dimensions. Here its Fourier coordinates are indexed by frequency,
$H$, labels, parts and entry coordinates, so the restricted dimension is
at most
\begin{equation}\label{op:orbit-dimension}
         C|O|\,|H|\,\dim_J U\le C|H|^4=\poly(r),
\end{equation}
where $C$ depends only on the fixed labels, and
$\dim_JU\le |H|^2$.  Form its entries using the explicit pulls and
Walsh signs, solve the two successive axis systems by Gaussian
elimination, and sum the output modes at $x$.  Together with
\eqref{op:orbit-weight}, this computes every local entry of
\[
                 E_{\mathrm{out}}((P_{\tau+1}-P_\tau)v).
\]
Matrix assembly, elimination and field arithmetic take $\poly(r)$
bit operations.  Canonical parsing and fixed defaults on malformed
values, illegal clocks or failed pivots make the local algorithm total;
on the promised data the prescribed matrices are invertible.

Use two cutoff-program buffers and one orbit-focus buffer.
Initially the target and all three buffers are zero.
At iteration $\tau$ keep the completed buffer for cutoff $\tau$
when $\tau>0$, and compare the target with it; at $\tau=0$ compare
the entire target with zero.  In the other, empty buffer fill the
cutoff-$\tau+1$ program in dependency order, and fill the orbit focus
in dependency order.  With both complete, perform one pivot adding
the local encoded increment to the target.  Its required comparison
now uses the completed cutoff-$\tau+1$ buffer.  Erase focus in reverse
order and erase the old cutoff buffer in reverse order when present.
Swap the buffer roles, determined by the iteration parity.
At cutoff $M_o$, compare the target directly with $X$ while erasing
the final cutoff buffer.  This finishes with the full desired target
and zero temporary scratch.

At every clock value enforce the full-code and duplication constraints
of the target, all initialized-prefix recurrences, inactive zeros,
and the stated comparison.  While a new proof is being prepared, a
completed old proof remains available.  At a pivot both the next
proof and the increment focus are complete.  In reverse traversal,
erased prefixes are first reconstructed in forward dependency order,
so the reverse pivot subtracts the same locally available increment.
Thus both directions are local, mutually inverse on canonical data,
and never use a partial certificate as a completed one.

For the size bound, each cutoff program uses at most
$O(|H|r^3)$ elementary stages up to fixed padding factors.
The orbit selectors and all right-$H$ paths use at most
$O(|H|r^2+|H|^3)$ such stages, with a polynomial allowance for
metadata.  Even a bound $O(r^{30})$ for both programs and their
descriptions suffices.  Each entry alphabet has at most $r^{18}$
field coordinates, each of $O(r)$ bits.  The two reused buffers
therefore have $\poly(r)$ bits per index.  There are $M_o$ iterations
of a common polynomial sublength, with a special initial omission
and final erasure; their location is decoded by arithmetic on a
logarithmic microclock.  No local step loops over all $M_o$ modes.
The different cutoffs reuse the same zeroed buffers and do not
contribute a factor $M_o$ to the word length.  Separate buffers for
the finitely many required A types keep every slot's type fixed.
These observations prove all bounds and claims of the lemma.
\end{proof}

\subsection{From sparse tape wiring to row systems}

We next construct a persistent tuple $P(b)$ independent of the clock
frame, where $b$ comprises the four tape-content and occupation arrays.
Its canonical formulas are defined for every $b\in J^{4h}$; Boolean
and phase-validity tests are additional conditions.
Testing a uniformly chosen tape row would not detect one false row with
constant probability. Instead we interpolate the row data, certify the
required products, and encode the remainder of each row-equation
polynomial modulo $Z^h-1$. A single false row makes that remainder
nonzero; comparing its full encoding with zero then has the code's
constant gap. The tuple retains the intermediate linear computations
needed to perform these checks at every frame.

\begin{lemma}[Persistent wiring and certification]\label{lem:outer-wiring}
There are $\poly(r)$ persistent full slots $P(b)$, with polynomial
canonical payload formulas in $b\in J^{4h}$, and packed local tests
using $\Pi$, with the following properties.
\begin{enumerate}[label=(\roman*)]
\item Structural tests uniquely determine $P(b)$ from its four logical
arrays.  Each variable factor array is a fixed permutation of one
logical array; a one-address logical edit changes that factor at at
most one publicly determined address.
\item Each primary has assigned logical degree at most a fixed $d_2$.
Every other persistent slot has a homogeneous linear recipe on the
primaries, including all intermediate matrix reencodings and all
phase-specific computations.  Its assigned logical degree can be
the maximum degree of its primary inputs.
\item There are persistent full encodings of
\begin{equation}\label{op:address-transform}
       U_D(x)=\sum_{j=0}^{h-1}(-1)^{x\cdot j}D_j
\end{equation}
for every needed logical singleton $D$, where $j$ denotes its padded
$r$-bit vector.  Its Walsh coefficient at $j<h$ is exactly $D_j$.
\item Conditional on exact duplicated full codes and a legal repeated
clock, either all structural formulas hold and $b$ is a valid
Boolean macrovertex at that clock, or one packed test rejects on
at least a fixed fraction $\gamma_{\mathrm{per}}>0$ of centers.
The same alternative without the macro-validity conclusion holds
if only structural tests are imposed.
\end{enumerate}
The canonical tuple and public data are deterministically constructible
in $\poly(N)$ time.  The number of field coordinates, public local
parameters and bit operations at one index is $\poly(r)$.
\end{lemma}

\begin{proof}
\paragraph{Sparse row equations.}
Use the normalization in Lemma~\ref{lem:macrograph}.
The Boolean tests for the four arrays are $z(z-1)=0$ at each row,
and content must vanish where occupation is zero, by $a(1-e)=0$.
For occupation in a prescribed order $\pi$, impose
\[
        e_{\pi(i+1)}(1-e_{\pi(i)})=0\quad(0\le i<h-1).
\]
Together with Booleanity, these equations say that occupied positions
form a prefix.  Its desired length is imposed separately by its
integer count.  Full occupation may instead be imposed row by row.
Fixed gate truth tables have multilinear representing polynomials
over $J$: for a Boolean input $a\in\{0,1\}^m$ its indicator is
$\prod_{i:a_i=1}z_i\prod_{i:a_i=0}(1-z_i)$, so summing the indicators
weighted by the truth-table values gives the desired polynomial.
The arity is fixed.  Gate, input-copy, designated-comparison, dummy,
and tape-to-tape-copy equations are conditioned by the occupation
of their destination where appropriate.  This expresses exactly
the phase rules in Section~\ref{sec:macrograph}; the phase prescription
ensures that an occupied computation only depends on earlier occupied
positions or the fixed full tape.  In a full $Q_x$, its $S$ inputs
are free, with the terminal and other-input relations imposed as in
its definition.  At shared phase boundaries fix either equivalent
description.

For each of the four phases and each unconditional family, organize
these equations into a fixed number of length-$h$ lists, padding with
zero equations.  Every operand map from list rows to an original
logical array has bounded multiplicity: this follows from bounded
fan-out for internal feeds, separate designated copies, and partial
permutations for adjacency and same-position copies.
An operand map of multiplicity at most $B$ splits into $B$ partial
injections by coloring its preimages separately at each target.
Each partial injection extends to a permutation of $[h]$, by any
bijection between the unused domain and range.  Separate logical
array types and intersect the finitely many operand-color and
active-row cases.  Because arity and $B$ are constants, this causes
only a constant number of cases.  Expand the representing polynomials
into monomials.  Each row list is now a sum of a fixed number of
bounded-length monomials, whose factors are permuted logical arrays
or public coefficient and mask arrays.  Empty monomials represent
constant terms.  Counting coefficient factors in the length still
gives a fixed bound on the degree.

\paragraph{Routing arbitrary permutations.}
For completeness a permutation of $2^\ell$ locations has a binary
switch realization with $2\ell-1$ layers, each layer switching pairs
along one fixed XOR bit, with a public Boolean mask constant on each
pair. This is the binary specialization of the Bene\v{s}
rearrangeability construction~\cite[Lemma~3 and Theorems~2--3]{Benes1964Permutation}.
Pair sources along the top
bit and do the same for targets.  The prescribed item routes form a
bipartite multigraph on these pairs with degree two at every vertex.
Every component is an even cycle, including a pair of parallel edges;
alternate its edges with two colors.  An initial top-bit switch and
a final top-bit switch route each item through the half corresponding
to its color.  Each half now has a permutation of $2^{\ell-1}$ items.
Recurse with the same layer schedule in the two halves.  The base case
$\ell=1$ has one layer and $\ell=0$ none; the recurrence for depth is
$L(\ell)=L(\ell-1)+2$.  Cycle coloring and recursive routing take
polynomial time in $h$.

Introduce a variable factor type for every routed intermediate, and
for each layer introduce the XOR-pulled counterpart $D^*$ of its input
$D$.  The layer output $D'$ satisfies the row equation
\begin{equation}\label{op:switch-row}
                      D'-D-m_{\mathrm{sw}}(D^*-D)=0.
\end{equation}
The mask is pair-constant, so this is exactly either the identity or
a transposition on each pair.  The counterpart identity itself is
checked by a scalar program below.  These equations are always
imposed, including in inactive phases.  Induction through the network
makes each variable factor a known permutation of an original logical
array, for arbitrary field values of that array.
There are only $O(r)$ layers among a fixed number of routing networks,
hence $\poly(r)$ factor and equation types with bounded arity per row.
For an edit of an original logical array at $j$, every factor derived
from it changes only at the image of $j$ under its fixed permutation.
All such addresses for one frame occupy $\poly(r)$ bits.
Their entire frame-indexed tables need not have a succinct random-access
formula; they will be stored as public encoded tables and focused
before use in Section~\ref{sec:outer-representation}.

\paragraph{Primaries and assigned degree.}
For every factor array $D\in J^h$, define its unique degree-$<h$
interpolant on the subgroup grid by
\begin{equation}\label{op:factor-interpolant}
                  F_D(Z)=\sum_{j=0}^{h-1}D_jL_j(Z).
\end{equation}
For each factor-occurrence list of each monomial, include primaries
for every subset $S$ of its occurrences, with payload the
length-$M_o$ coefficient table, padded by zeros, of
\begin{equation}\label{op:subset-primary}
                         F_S(Z)=\prod_{a\in S}F_{D_a}(Z).
\end{equation}
The index set here consists of occurrences: repeated uses of one
factor remain distinct.  The empty subset gives $1$.
Identify shared singleton primaries or impose direct full-word
equalities on copies, and include independent singletons when needed.
The assigned logical degree $k_S$ counts variable occurrences only;
public factors have degree zero.  Let $d_2$ bound the length of every
factor-occurrence list after all switch and row equations have been
included.  It is a fixed constant, and $k_S\le d_2$.
The honest ordinary polynomial degree is at most $d_2(h-1)$.
Choose the earlier size constant so that
\begin{equation}\label{op:degree-slack}
                         M_o>(8d_2+8)h.
\end{equation}
All primary formulas exist for arbitrary $b\in J^{4h}$.

\paragraph{Index independence and support.}
Scalar primaries are intended to depend only on $x$, not on $u$.
Test invariance of their payloads under left $(0,a)$ pulls for the
fixed generating list of $H$.  Such a pull fixes $x$ and changes
$u$ to $au$, and shares B with the original untwisted scalar slot.
Compare whole systematic A rows as above, using the pulled local
expression without introducing a free encoder witness.
The shared-B exact gap forces invariance under each generator;
since these generate $H$, this is precisely independence from $u$.
No independent invariance test is needed on a derived program stage
whose exact recipe already determines it, and matrix intermediates
need not be $u$-independent.

Use range masks from Lemma~\ref{op:scalar-programs} to require
singleton coefficient support in $[0,h)$ and product-primary support
in $[0,d_2h+1)$, with sharper public bounds when appropriate.
Encode the masked complement as a full output and compare that
entire word with zero.  A nonzero output has full-code distance.
Compare each public singleton and each empty product with its
precomputed canonical full encoding.  These are structural conditions.

\paragraph{Product certification.}
Compute the length-$M_o$ DFT of each relevant primary coefficient
table by certified full stages.  Its payload at integer $y$ is the
polynomial's value at $\zeta^y$.  For every nontrivial subset choose
one occurrence $a\in S$ and check, at one fixed systematic payload
cell at $s=t=0$,
\begin{equation}\label{op:primary-product-test}
                   F_S(\zeta^y)
                  =F_{S\setminus\{a\}}(\zeta^y)F_{D_a}(\zeta^y).
\end{equation}
These are pointwise nonlinear tests; the stored transform stages
themselves have linear recipes.  Once support, independence and
transform identities hold, a false equality is the nonzero polynomial
$F_S-F_{S\setminus\{a\}}F_{D_a}$, of degree less than
$(d_2+1)h+1<M_o/2$.  A nonzero polynomial of degree $D$ over a field
has at most $D$ roots, as follows by successively dividing by $Z-z$
at distinct roots.  Thus the false equality fails at more than half
the $M_o$ distinct points $\zeta^y$.  Uniform $k$ has uniform $x$,
and independence from $u$ prevents the extra index from diluting this
gap.  Induction on subset size forces exactly all products
\eqref{op:subset-primary} once these tests hold.

\paragraph{Counterparts and the subgroup annihilator.}
Apply the length-$h$ Fourier transform to the low $\ell$ bits of each
singleton's padded coefficient table.  Its first block is the value
table $D$, and every higher block is zero by coefficient support.
For each XOR counterpart use the certified bit pull on this padded
table and compare its full output with the counterpart's separately
computed padded value table.  Their full-operation gap forces the
exact counterpart identity needed in \eqref{op:switch-row}.

For a row-equation list let $R(Z)$ be the indicated linear combination
of its product primaries.  Its value at $\omega^j$ is exactly the
row expression at $j$.  Its remainder modulo $Z^h-1$ has coefficients
\begin{equation}\label{op:remainder-fold}
       [Z^i](R\bmod(Z^h-1))
          =\sum_{z=0}^{M_o/h-1}[Z^{i+zh}]R,
                \quad 0\le i<h,
\end{equation}
and zero padding elsewhere.  This formula follows from
$Z^{i+zh}\equiv Z^i\pmod{Z^h-1}$, since the coefficient table is
supported below $M_o$.  It is a short linear program: apply
$\prod_{i=\ell}^{r-1}(I+R_{(e_i,1)})$ to sum all high-bit blocks,
then mask to high bits zero.  Require its entire full word to be
zero when the equation is imposed.

The polynomial $Z^h-1$ is exactly the annihilator of the subgroup grid:
it has the $h$ distinct roots $\omega^j$, because $h$ is invertible
in $J$, and its degree is $h$.  The remainder has degree below $h$.
Therefore it vanishes as a polynomial if and only if $R$ vanishes at
all grid points.  In one direction divisibility suffices; in the
other direction a degree-$<h$ polynomial with $h$ distinct roots
is zero.  This proves that the folded full-word checks impose the
row equations exactly, without any random sampling loss over rows.
All structural switch equations and their folded computations are
always present.  All phase-specific folded computations are also
stored persistently; the repeated clock selects only which final
zero conditions are imposed.

\paragraph{Counts and address transforms.}
The constant coefficient of \eqref{op:factor-interpolant} is
\begin{equation}\label{op:occupation-count}
                       [Z^0]F_e=h^{-1}\sum_{j=0}^{h-1}e_j,
\end{equation}
since every $L_j$ has constant coefficient $h^{-1}$.
Mask the occupation coefficient table to index zero, using all bit
masks.  Compare its full word with the public full encoding of
$1[\underline x=0]$ times $c h^{-1}$, where the expected count
$c\in[0,h]$ is computed from the legal repeated stage and phase.
This scalar is used in the comparison; it is not part of a changing
persistent value.  On Boolean occupation data, characteristic
$p>M_o\ge h$ implies that two counts in $[0,h]$ are equal in $J$
only when they are equal as integers.  Together with monotonicity
this enforces the prescribed prefix exactly.

For address access, take the padded length-$h$ value table obtained
above and apply the full $r$-bit Walsh transform.  Its result is
\eqref{op:address-transform}.  Walsh orthogonality gives
\[
 M_o^{-1}\sum_x U_D(x)(-1)^{x\cdot j}
 =\sum_{i<h}D_i
        \left(M_o^{-1}\sum_x(-1)^{x\cdot(i+j)}\right)
 =D_j \qquad(j<h).
\]
In particular the high bits of the zero-padded value table cause
$U_D$ to repeat across the corresponding high output blocks;
they do not introduce an additional normalization factor.

\paragraph{Exact alternatives and complexity.}
Every stage of every mask, transform, counterpart comparison and
remainder computation is stored as a full slot with its specified
linear recipe, including mixed and untwisted matrix intermediates.
Give a derived stage the maximum assigned degree of its primary
ancestors and retain its ancestor subprogram or shared directed
acyclic recipe.  Public indexed arrays are degree-zero inputs with
precomputed encoded constants.

On exact duplicated full codewords, a false independence, support,
public-constant, full-operation, counterpart, folded-zero or count
identity has a constant gap by the shared-axis or full-code distance.
If none is false, the polynomial root argument forces all primary
products.  Singleton support makes the interpolant unique from its
subgroup values.  The counterpart and switch equations then determine
all factors as their prescribed permutations of the four original
logical arrays, and the public factors are fixed.  Thus the structural
tests force precisely $P(b)$.  The remaining Boolean, prefix, count
and selected-phase row equations are exactly the macrovertex rules.
If any is false, its folded or count check has the same exact gap.
All tests are packed at the center; choosing one false relation in
this argument entails no division by their number.  The minimum of
the fixed code gaps and the product gap gives
$\gamma_{\mathrm{per}}>0$ and proves both alternatives.

There are $O(r)$ routed factor and row families before the scalar
program expansions, a constant number of occurrence subsets per
monomial, and polynomially many stages per transform.  Thus the
number of full slots, recipes and local metadata is $\poly(r)$.
The entry dimension and field bit length give a further fixed
polynomial factor.  At one center the checker scans these lists,
performs field operations and evaluates fixed-length control
decisions, all in $\poly(r)$ bit operations.  It reads public encoded
constants only at the publicly specified positions in $\Pi(k)$;
arbitrary hidden-frame table lookup is never assumed.
All public permutations, masks, coefficients, encodings and complete
canonical arrays are constructible globally by polynomial-time
combinatorial and linear-algebra algorithms.  Since
$|\Gamma|=M_o|H|=O(N\operatorname{polylog} N)$ and every scalarized encoding system
has polynomial dimension in $N$, this takes $\poly(N)$ time.
The persistent word length is $O(N\operatorname{polylog} N)$ field-bit storage,
with a fixed exponent in the polylogarithmic factor.  This completes
the construction and proof.
\end{proof}
